%%%%%%%%%%%%%%%%%%%%%%%%%%%%%%%%%%%%%%%%%%%%%%%%%%%%%%%%%%%%%%%%%%%%%
%%%%%%%%%%%%%%%%%%%%%%%%% Resumo %%%%%%%%%%%%%%%%%%%%%%%%%%%%%%%%%%%%
%%%%%%%%%%%%%%%%%%%%%%%%%%%%%%%%%%%%%%%%%%%%%%%%%%%%%%%%%%%%%%%%%%%%%

%====================================================================
% RESUMO
\setlength{\absparsep}{18pt} % ajusta o espaçamento dos parágrafos do resumo
\begin{resumo}

Indicações dos pontos mais relevantes em um texto, passando uma rápida e clara visualização dos conteúdos contidos no trabalho. Deverá possuir palavras-chaves logo abaixo da última linha de texto. De preferência utiliza-se aproximadamente três palavras-chaves para representar o conteúdo do trabalho. Devendo estar em português, separadas por ponto, com a primeira letra em maiúsculo e o restante em minúsculo.

 \textbf{Palavras-chave}: modelo de artigo. resumo. template.   
\end{resumo}

%==================================================================